% archon:covers AlgebraicJacobian/RiemannRoch/OCofP.lean

\chapter{The line bundle \(\mathcal O_C(P)\) and its global sections (RR.3)}
\label{chap:RiemannRoch_OCofP}

% SOURCE: Hartshorne, \emph{Algebraic Geometry}, II.6, p.~144 (definition
% of \(\mathscr L(D)\)); II.7, Proposition~7.7, p.~157 (effective divisors
% as divisors of zeros of global sections); IV.1, pp.~294--297
% (Riemann--Roch, Lemma~1.2, Examples~1.3.3--1.3.5)
% (read from references/hartshorne-algebraic-geometry.pdf, PDF pages 161,
% 174, 311--314). The verbatim quotes pasted in the declaration blocks
% below were copied character-by-character from the rendered page images
% of the scanned PDF (no text layer is present).

\section{Setup and motivation}

This chapter is the third of four sub-build chapters \texttt{RR.1}--\texttt{RR.4}
that close the project's headline \emph{RR bridge}
\[
  \mathtt{genusZero\_curve\_iso\_P1}
  \,\colon\,
  \text{a smooth proper geometrically irreducible curve \(C / \bar k\) with
  \(g(C) = 0\) is isomorphic to \(\mathbb P^1_{\bar k}\).}
\]
The bridge as proved in Hartshorne (Example~IV.1.3.5) chains together:
the Weil-divisor group \(\mathrm{Div}(C)\), the degree map $\deg \colon \mathrm{Div}(C)
\to \mathbb Z$, and the divisor $[P]$ of a closed point
(\texttt{RR.1}, \cref{def:divisor_closed_point}); the Riemann--Roch
dimension formula in genus~\(0\)
\[
  \ell(D) \;=\; \deg(D) + 1 \qquad \text{for } \deg(D) \geq 0
\]
via the Euler-characteristic identity
(\texttt{RR.2}, \cref{thm:euler_char_eq_deg_plus_one_minus_genus}); the global-sections
calculation \(\dim_{\bar k} H^0(C, \mathcal O_C(P)) = 2\) on a smooth
proper genus-\(0\) curve (this chapter, \texttt{RR.3}); the
``rational curve \(\Rightarrow \cong \mathbb P^1\)'' classification via
\texttt{Proj.fromOfGlobalSections} (\texttt{RR.4}).

The contribution of \texttt{RR.3} is twofold. First, it constructs the
invertible sheaf \(\mathcal O_C(P)\) associated to a closed point \(P \in C\)
on a smooth proper curve, identifies its \(\bar k\)-vector space of global
sections with the \emph{Riemann--Roch space}
\[
  L([P]) \;=\;
  \bigl\{f \in K(C)^{\times} \;\big|\; \div(f) + [P] \geq 0\bigr\}
  \,\cup\, \{0\},
\]
and computes its dimension as \(\dim_{\bar k} H^0(C, \mathcal O_C(P)) = 2\)
in the genus-\(0\) specialisation. Second, it deduces from the
two-dimensionality that \(C\) admits a non-constant rational function \(f\)
having at most a simple pole at \(P\) and no other poles. This rational
function is the seed for \texttt{RR.4}: the pair $(1, f) \in
H^0(C, \mathcal O_C(P))^{\oplus 2}$ generates the sheaf and, via
\texttt{Proj.fromOfGlobalSections}, produces the morphism
\(C \to \mathbb P^1_{\bar k}\) whose degree-\(1\) identification is the
content of \texttt{RR.4}.

% NOTE (Direct cohomological vanishing vs. Serre duality — iter-174
% sub-phase choice). The dimension computation
% \(\dim_{\bar k} H^0(C, \mathcal O_C(P)) = 2\) would, in Hartshorne's
% original presentation, route through Serre duality
% \(H^1(C, \mathcal L) \cong H^0(C, \omega_C \otimes \mathcal L^{-1})^*\)
% to vanish \(H^1(C, \mathcal O_C(P))\) from the degree count
% \(\deg(K_C - [P]) = 2g - 2 - 1 = -3 < 0\). Mathlib (snapshot
% \texttt{b80f227}) ships no dualizing sheaf \(\omega_X\) on a scheme and
% no Serre duality theorem. The genus-\(0\) specialisation does \emph{not}
% require Serre duality: the vanishing \(H^1(C, \mathcal O_C(P)) = 0\) is
% delivered by the long exact sequence of the structural short exact
% sequence
% \[
%   0 \to \mathcal O_C \to \mathcal O_C(P) \to k(P) \to 0
% \]
% together with the genus-\(0\) hypothesis
% \(\dim_{\bar k} H^1(C, \mathcal O_C) = g = 0\) and the skyscraper
% identity \(H^1(C, k(P)) = 0\). The chapter therefore uses the direct
% cohomological vanishing route exclusively; Serre duality is left for
% the post-RR-bridge follow-up chapter that handles the general-\(g\)
% statement.

\paragraph{Standing hypotheses.} Throughout this chapter, \(\bar k\) is an
algebraically closed field and \(C\) is a smooth proper geometrically
irreducible curve over \(\bar k\). Concretely, the Lean signatures
(\texttt{AlgebraicJacobian/RiemannRoch/OCofP.lean}) carry the same
typeclass package as the rest of the project's curve API:
\begin{quote}
  \texttt{\{kbar : Type u\} [Field kbar] [IsAlgClosed kbar]\\
   (C : Over (Spec (.of kbar))) [IsProper C.hom]\\
   \phantom{}\quad[SmoothOfRelativeDimension 1 C.hom]\\
   \phantom{}\quad[GeometricallyIrreducible C.hom] [IsIntegral C.left]}
\end{quote}
The genus \(g(C)\) is the one defined in \cref{def:genus} of
\texttt{Genus.tex}; the divisor \([P]\) of a closed point is the one of
\cref{def:divisor_closed_point} of \texttt{RiemannRoch\_WeilDivisor.tex};
the principal divisor \(\div(f)\) of a nonzero rational function is the one
of \cref{def:principal_divisor}, and the order \(\ord_Q(f)\) at a prime
divisor is the standard discrete valuation of the project's Weil-divisor
order API.

\section{The line bundle of a closed point on a curve}

% --- iter-187 plan-phase additions: pins for the iter-187 refactor's
% Step 3 + Step 4 sub-helpers (carrierPresheaf functor + sheaf property)
% which are the direct-construction substrate the definition below
% consumes (via the iter-186 plan-phase `analogies/ocofp-carrierset-submodule-api.md`
% Decision 3 + Decision 4 recipe). The Lean declarations are produced
% by the iter-187 refactor-`ocofp-steps3to5` dispatch. The substantive
% mathematical content is in the Definition block below; these helper
% blocks pin the substrate hooks the prover-side construction uses.

\begin{definition}\leanok
[Carrier set of the line bundle of a closed point over an open]
  \label{def:lineBundleAtClosedPoint_carrierSet}
  \lean{AlgebraicGeometry.Scheme.lineBundleAtClosedPoint.carrierSet}
  Let \(C\) be an integral, locally Noetherian, regular-in-codimension-one
  scheme over \(\bar k\), let \(P \in C\) be a point of coheight \(1\), and
  let \(U \subseteq C\) be an open set. The \emph{carrier set}
  \(\mathrm{carrierSet}_P(U) \subseteq K(C)\) is the set of rational
  functions \(f\) satisfying the order conditions \(\ord_Q(f) \ge 0\) for
  every prime divisor \(Q\) with \(Q \in U\) and \(Q \ne P\), together with
  \(\ord_P(f) \ge -1\) whenever \(P \in U\). It is the underlying set of the
  section module of \(\mathcal O_C(P)\) over \(U\).
\end{definition}

\begin{lemma}\leanok
[Monotonicity of the carrier set]
  \label{lem:lineBundleAtClosedPoint_carrierSet_mono}
  \lean{AlgebraicGeometry.Scheme.lineBundleAtClosedPoint.carrierSet_mono}
  \uses{def:lineBundleAtClosedPoint_carrierSet}
  With \(C, P\) as in \cref{def:lineBundleAtClosedPoint_carrierSet}, if
  \(V \subseteq U\) are opens of \(C\) then
  \(\mathrm{carrierSet}_P(U) \subseteq \mathrm{carrierSet}_P(V)\): the order
  conditions over the smaller open \(V\) involve fewer prime divisors and are
  therefore implied by those over \(U\). This antitone inclusion is the
  restriction map of the carrier presheaf.
\end{lemma}

\begin{proof} Proved directly in Lean. \end{proof}

\begin{definition}\leanok
[Nonemptiness of the top open of an integral scheme]
  \label{def:lineBundleAtClosedPoint_instNonemptyTopOpen}
  \lean{AlgebraicGeometry.Scheme.lineBundleAtClosedPoint.instNonemptyTopOpen}
  For an integral scheme \(X\), the top open \(\top \subseteq X\) is nonempty
  as a scheme: its underlying set is all of \(X\), and \(X\) is nonempty since
  an integral scheme is irreducible. This instance lets the standard algebra
  map \(\Gamma(X, U) \to K(X)\) fire at \(U = \top\).
\end{definition}

\begin{definition}\leanok
[The \(\bar k\)-algebra structure on the function field]
  \label{def:lineBundleAtClosedPoint_instAlgebraKbarFunctionField}
  \lean{AlgebraicGeometry.Scheme.lineBundleAtClosedPoint.instAlgebraKbarFunctionField}
  \uses{def:lineBundleAtClosedPoint_instNonemptyTopOpen}
  For an integral scheme \(C\) over \(\bar k\), the function field \(K(C)\)
  carries a \(\bar k\)-algebra structure, obtained by composing the algebra
  map \(\bar k \to \Gamma(C, \top)\) with the canonical map
  \(\Gamma(C, \top) \to K(C)\). This structure is what makes \(K(C)\) a
  \(\bar k\)-module, so that the carrier sets can be upgraded to
  \(\bar k\)-submodules.
\end{definition}

\begin{definition}\leanok
[Carrier submodule of the line bundle of a closed point]
  \label{def:lineBundleAtClosedPoint_carrierSubmodule}
  \lean{AlgebraicGeometry.Scheme.lineBundleAtClosedPoint.carrierSubmodule}
  \uses{def:lineBundleAtClosedPoint_carrierSet,
        def:lineBundleAtClosedPoint_instAlgebraKbarFunctionField}
  With \(C, P, U\) as in \cref{def:lineBundleAtClosedPoint_carrierSet}, the
  carrier set \(\mathrm{carrierSet}_P(U)\) is a \(\bar k\)-submodule of
  \(K(C)\). It contains \(0\); it is closed under addition by the
  non-archimedean inequality \(\ord_Q(f + g) \ge \min(\ord_Q f, \ord_Q g)\)
  for the discrete valuation at each regular codimension-one stalk; and it is
  closed under scaling by a constant \(c \in \bar k\), since \(c\) maps to a
  unit in every stalk and hence has order \(0\) at every prime divisor.
\end{definition}

\begin{definition}\leanok
[Bot-trivialisation submodule]
  \label{def:lineBundleAtClosedPoint_trivAtBot}
  \lean{AlgebraicGeometry.Scheme.lineBundleAtClosedPoint.trivAtBot}
  \uses{def:lineBundleAtClosedPoint_instAlgebraKbarFunctionField}
  For an integral scheme \(C\) over \(\bar k\) and an open \(U \subseteq C\),
  the \(\bar k\)-submodule \(\mathrm{trivAtBot}(U) \subseteq K(C)\) equals all
  of \(K(C)\) when \(U \ne \emptyset\) and equals \(\{0\}\) when
  \(U = \emptyset\). Intersecting the carrier submodule with this factor
  enforces the sheaf-at-\(\emptyset\) condition \(F(\emptyset) = 0\) required
  by the Grothendieck topology of opens.
\end{definition}

\begin{lemma}\leanok
[Monotonicity of the sheaf-corrected carrier submodule]
  \label{lem:lineBundleAtClosedPoint_carrierSubmoduleSheaf_le}
  \lean{AlgebraicGeometry.Scheme.lineBundleAtClosedPoint.carrierSubmoduleSheaf_le}
  \uses{def:lineBundleAtClosedPoint_carrierSubmoduleSheaf,
        lem:lineBundleAtClosedPoint_carrierSet_mono}
  With \(C, P\) as above, if \(V \subseteq U\) are opens with
  \(V \ne \emptyset\), then
  \(\mathrm{carrierSubmoduleSheaf}_P(U) \le \mathrm{carrierSubmoduleSheaf}_P(V)\).
  The hypothesis \(V \ne \emptyset\) is necessary, since a nonzero section
  over \(U\) need not satisfy the trivialisation-at-\(\emptyset\) condition.
  This inclusion supplies the restriction maps of the carrier presheaf on
  nonempty opens.
\end{lemma}

\begin{proof}
  \uses{lem:lineBundleAtClosedPoint_carrierSet_mono}
  Proved directly in Lean, from \cref{lem:lineBundleAtClosedPoint_carrierSet_mono}.
\end{proof}

\begin{definition}\leanok
[Type-level subfunctor presentation of the carrier]
  \label{def:lineBundleAtClosedPoint_carrierTypeSubfunctor}
  \lean{AlgebraicGeometry.Scheme.lineBundleAtClosedPoint.carrierTypeSubfunctor}
  \uses{def:lineBundleAtClosedPoint_carrierSubmoduleSheaf}
  With \(C, P\) as above, the carrier of \(\mathcal O_C(P)\) is presented as a
  subfunctor of the type-valued presheaf \(U \mapsto (U \to K(C))\) of
  arbitrary \(K(C)\)-valued functions: a section over \(U\) is a constant
  function with value \(f \in \mathrm{carrierSubmoduleSheaf}_P(U)\). This
  presentation reduces the sheaf condition to a stalk-locality check against
  the ambient sheaf of type-valued functions, exploiting that any two nonempty
  opens of the irreducible scheme \(C\) intersect.
\end{definition}

\begin{definition}

\leanok
[Presheaf carrier for the line bundle of a closed point]
  \label{def:lineBundleAtClosedPoint_carrierPresheaf}
  \lean{AlgebraicGeometry.Scheme.lineBundleAtClosedPoint.carrierPresheaf}
  \uses{def:lineBundleAtClosedPoint_carrierSubmoduleSheaf,
        lem:lineBundleAtClosedPoint_carrierSubmoduleSheaf_le}
  For each open \(U \subseteq C\), the \(\bar k\)-submodule
  \(\mathrm{carrierSubmodule}_P(U) \subset K(C)\) of section candidates
  satisfying the order conditions ``\(\mathrm{ord}_Q f \ge 0\) for all
  prime divisors \(Q \in U\) other than \(P\), and additionally
  \(\mathrm{ord}_P f \ge -1\) when \(P \in U\)'' assembles into a
  \(\bar k\)-linear functor
  \[
    \texttt{carrierPresheaf}_P : (\Opens\, |C|)^{\op}
    \longrightarrow \ModuleCat_{\bar k},
    \qquad U \longmapsto \mathrm{carrierSubmodule}_P(U),
  \]
  with restriction \(V \subseteq U \mapsto \texttt{Submodule.inclusion}\,
  (\texttt{carrierSet\_mono})\) on each pair. The functor laws follow from
  the antitone monotonicity \(U \mapsto \mathrm{carrierSet}_P(U)\) (witness
  \(\texttt{carrierSet\_mono}\) at L181-192 of the Lean file) by
  \(\texttt{ext}\,\langle x, \_\rangle\); \texttt{rfl}.
\end{definition}

\begin{lemma}

\leanok
[Sheaf property of the carrier presheaf]
  \label{lem:lineBundleAtClosedPoint_carrierPresheaf_isSheaf}
  \lean{AlgebraicGeometry.Scheme.lineBundleAtClosedPoint.carrierPresheaf_isSheaf}
  \uses{def:lineBundleAtClosedPoint_carrierPresheaf,
        def:lineBundleAtClosedPoint_carrierSubmoduleSheaf,
        def:lineBundleAtClosedPoint_carrierTypeSubfunctor,
        def:lineBundleAtClosedPoint_trivAtBot}
  The presheaf \texttt{carrierPresheaf}\(_P\) (defined via the
  carrier-refinement of \cref{def:lineBundleAtClosedPoint_carrierSubmoduleSheaf})
  is a sheaf in the appropriate Grothendieck topology on \(C\). Concretely,
  the locality + gluing conditions for \(\ModuleCat_{\bar k}\)-valued sections
  reduce (via \(\texttt{isSheaf\_iff\_isSheaf\_forget}\)) to the sheaf
  condition on the underlying Set-valued functor: each
  \(\mathrm{carrierSubmoduleSheaf}_P(U)\) is a subset of the constant
  \(K(C)\)-sheaf restricted to \(U\), and the sheaf property in the
  constant-sheaf base case is direct. The project's
  \texttt{toModuleKSheaf} template
  (\texttt{Cohomology/StructureSheafModuleK/SheafProperty.lean}) is
  the model for this argument.
\end{lemma}

\begin{definition}


\leanok
[Carrier-submodule refinement enforcing \(F(\bot)=0\)]
  \label{def:lineBundleAtClosedPoint_carrierSubmoduleSheaf}
  \lean{AlgebraicGeometry.Scheme.lineBundleAtClosedPoint.carrierSubmoduleSheaf}
  \uses{def:lineBundleAtClosedPoint_carrierSubmodule,
        def:lineBundleAtClosedPoint_trivAtBot}
  Iter-188 Lane A discovery: the raw \(\mathrm{carrierSubmodule}_P\) of
  \cref{def:lineBundleAtClosedPoint_carrierPresheaf} satisfies
  \(\mathrm{carrierSubmodule}_P(\bot) = K(C)\) (since the order conditions
  are vacuously satisfied on the empty open). This violates the sheaf
  axiom on the Grothendieck topology of opens: the empty sieve covers
  \(\bot\), forcing \(F(\bot) =\) terminal \(= 0\) on
  \(\ModuleCat_{\bar k}\)-valued functors.
  The refinement
  \[
    \mathrm{carrierSubmoduleSheaf}_P(U) \;:=\;
    \mathrm{carrierSubmodule}_P(U) \;\sqcap\; \mathrm{trivAtBot}\,U,
    \qquad \mathrm{trivAtBot}\,U \;:=\;
    \{\, f \;\mid\; U \ne \bot \,\lor\, f = 0 \,\}.
  \]
  forces \(F(\bot) = \bot\) (intersection annihilates) while preserving
  \(F(U)\) on non-empty \(U\) (since then \(\mathrm{trivAtBot}\,U = \top\)).
  The case-based restriction map (zero target on \(\bot\); inclusion
  otherwise) gives a well-formed sheaf-condition substrate consumed by
  \cref{lem:lineBundleAtClosedPoint_carrierPresheaf_isSheaf}. Downstream
  consumers \texttt{lineBundleAtClosedPoint}, \texttt{toFunctionField},
  \texttt{globalSections\_iff\_*} are updated to supply the
  \texttt{trivAtBot} witness at \(\op\,\top\) via the \texttt{htop\_ne\_bot}
  helper.
\end{definition}

\begin{definition}



\leanok
[Line bundle of a closed point]
  \label{def:lineBundleAtClosedPoint}
  \lean{AlgebraicGeometry.Scheme.lineBundleAtClosedPoint}
  \uses{def:lineBundleAtClosedPoint_carrierPresheaf, lem:lineBundleAtClosedPoint_carrierPresheaf_isSheaf}
  
  % SOURCE: [Hartshorne], II.6, p.~144 (definition of the invertible
  % sheaf \(\mathscr L(D)\) associated to a Cartier divisor) (read from
  % references/hartshorne-algebraic-geometry.pdf, PDF page 161).
  % SOURCE QUOTE: "Definition. Let \(D\) be a Cartier divisor on a scheme
  % \(X\), represented by \(\{(U_i, f_i)\}\) as above. We define a subsheaf
  % \(\mathscr L(D)\) of the sheaf of total quotient rings \(\mathscr K\) by
  % taking \(\mathscr L(D)\) to be the sub-\(\mathscr O_X\)-module of
  % \(\mathscr K\) generated by \(f_i^{-1}\) on \(U_i\). This is well-defined,
  % since \(f_i/f_j\) is invertible on \(U_i \cap U_j\), and \(f_i^{-1}\) and
  % \(f_j^{-1}\) generate the same \(\mathscr O_X\)-module. We call
  % \(\mathscr L(D)\) the \emph{sheaf associated to \(D\)}."
  % SOURCE QUOTE: "Proposition 6.13. \emph{Let \(X\) be a scheme. Then:
  % (a) for any Cartier divisor \(D\), \(\mathscr L(D)\) is an invertible
  % sheaf on \(X\).}"
  \textit{Source: Hartshorne, II.6, p.~144 (definition of
  \(\mathscr L(D)\) and Proposition~6.13(a)).}
  Let \(C\) be a smooth proper curve over \(\bar k\) and let \(P \in C\) be a
  closed point. Let \(\mathfrak m_P \subset \mathcal O_{C, P}\) denote the
  maximal ideal of the local ring at \(P\); by smoothness of \(C\),
  \(\mathcal O_{C, P}\) is a discrete valuation ring and \(\mathfrak m_P\)
  is generated by any uniformizer \(f_P \in \mathcal O_{C, P}\) (an
  element of \(\mathfrak m_P \setminus \mathfrak m_P^2\)). The Cartier
  divisor associated to the closed point \(P\) is the effective
  Cartier divisor cut out locally by the uniformizer \(f_P\) on an open
  affine neighbourhood of \(P\), and by the unit \(1\) on the complement
  \(C \setminus \{P\}\). The \emph{line bundle of the closed point \(P\)},
  denoted \(\mathcal O_C(P)\), is the locally free \(\mathcal O_C\)-module
  of rank \(1\) defined by Hartshorne's construction \(\mathscr L(D)\)
  applied to this Cartier divisor: explicitly, \(\mathcal O_C(P)\) is the
  sub-\(\mathcal O_C\)-module of the sheaf of total quotient rings
  \(\mathcal K_C\) (the constant sheaf of the function field \(K(C)\) on the
  integral scheme \(C\)) that is generated by \(f_P^{-1} \in K(C)\) on a
  neighbourhood of \(P\) and by \(1 \in K(C)\) on the complement
  \(C \setminus \{P\}\). Hartshorne's Proposition~6.13(a) confirms that
  this subsheaf is an invertible sheaf, independently of the choice of
  uniformizer \(f_P\) (any two uniformizers differ by a unit in
  \(\mathcal O_{C, P}^{\times}\), so they generate the same
  \(\mathcal O_C\)-submodule of \(\mathcal K_C\)).

  \paragraph{Equivalent description via the ideal sheaf.} Let
  \(\mathcal I_P \subset \mathcal O_C\) denote the ideal sheaf of the
  closed point \(P\) (locally generated by the uniformizer \(f_P\) near
  \(P\), and equal to \(\mathcal O_C\) on the complement). Then
  \(\mathcal I_P\) is the line bundle \(\mathcal O_C(-P)\) associated to the
  Weil divisor \(-[P]\), and \(\mathcal O_C(P)\) is the
  \(\mathcal O_C\)-linear dual:
  \[
    \mathcal O_C(P) \;\cong\; \mathcal H\!om_{\mathcal O_C}(\mathcal I_P,
    \mathcal O_C)
    \;\cong\; \mathcal I_P^{-1}.
  \]
  Both descriptions yield the same invertible sheaf; the
  ``subsheaf of \(\mathcal K_C\)'' description above is the one used in
  the prose, while the ``dual of the ideal sheaf'' description is the
  most direct entry point for the Lean construction (Mathlib's
  \texttt{AlgebraicGeometry.Scheme.IdealSheafData} packages
  quasi-coherent ideal sheaves, from which the dual is constructed via
  the internal-hom on invertible sheaves).

  \paragraph{Lean signature scope.} The Lean definition
  \texttt{AlgebraicGeometry.Scheme.lineBundleAtClosedPoint} takes a
  scheme \(C\), a point \(P : C\), and a closedness witness
  \(\texttt{IsClosed (\{P\} : Set C)}\) (the same input shape as
  \cref{def:divisor_closed_point}), and returns an invertible
  \(\mathcal O_C\)-module \(\mathcal O_C(P)\). The well-definedness as an
  invertible sheaf is part of the construction and goes through
  smoothness of \(C\) (delivering the DVR stalk at \(P\)) and the standard
  ideal-sheaf-of-a-closed-point machinery in Mathlib's
  \texttt{AlgebraicGeometry.IdealSheaf} chapter. The chapter's
  consumption of the closed-point hypothesis matches \texttt{RR.1}:
  closed plus dimension-\(1\)-integral delivers \(\mathrm{codim}_C \{P\} = 1\)
  automatically. The declaration is \texttt{noncomputable}.
\end{definition}

\begin{definition}


\leanok
[Inclusion of \(\mathcal O_C(P)\) into the function field]
  \label{def:lineBundleAtClosedPoint_toFunctionField}
  \lean{AlgebraicGeometry.Scheme.lineBundleAtClosedPoint.toFunctionField}
  
  % SOURCE: [Hartshorne], II.6, p.~144 (definition of \(\mathscr L(D)\) as a
  % subsheaf of the sheaf of total quotient rings \(\mathscr K\); the
  % canonical inclusion \(\mathscr L(D) \hookrightarrow \mathscr K\) is part
  % of the construction) (read from
  % references/hartshorne-algebraic-geometry.pdf, PDF page 161).
  % SOURCE QUOTE: "Definition. Let \(D\) be a Cartier divisor on a scheme
  % \(X\), represented by \(\{(U_i, f_i)\}\) as above. We define a subsheaf
  % \(\mathscr L(D)\) of the sheaf of total quotient rings \(\mathscr K\) by
  % taking \(\mathscr L(D)\) to be the sub-\(\mathscr O_X\)-module of
  % \(\mathscr K\) generated by \(f_i^{-1}\) on \(U_i\). This is well-defined,
  % since \(f_i/f_j\) is invertible on \(U_i \cap U_j\), and \(f_i^{-1}\) and
  % \(f_j^{-1}\) generate the same \(\mathscr O_X\)-module. We call
  % \(\mathscr L(D)\) the \emph{sheaf associated to \(D\)}."
  \textit{Source: Hartshorne, II.6, p.~144 (definition of \(\mathscr L(D)\)
  as a subsheaf of \(\mathscr K\)).}
  Let \(C\) be a smooth proper geometrically irreducible curve over \(\bar k\)
  with \(C\) integral, and let \(P \in C\) be a closed point. The line bundle
  \(\mathcal O_C(P)\) of \cref{def:lineBundleAtClosedPoint} is by
  construction a sub-\(\mathcal O_C\)-module of the sheaf of total quotient
  rings \(\mathcal K_C\). Since \(C\) is integral, \(\mathcal K_C\) is the
  constant sheaf at the function field \(K(C)\), and taking global sections
  of the structural inclusion \(\mathcal O_C(P) \hookrightarrow \mathcal K_C\)
  produces a canonical \(\bar k\)-linear map
  \[
    \iota \;\colon\; H^0(C, \mathcal O_C(P))
      \;\longrightarrow\; H^0(C, \mathcal K_C) \;\cong\; K(C)
  \]
  sending a global section \(s \in H^0(C, \mathcal O_C(P))\) to the rational
  function on \(C\) obtained by composing \(s\) with the inclusion of
  \(\mathcal O_C(P)\) into \(\mathcal K_C\) at each point. The map \(\iota\) is
  injective (the sub-\(\mathcal O_C\)-module $\mathcal O_C(P) \subset
  \mathcal K_C$ is torsion-free of rank $1$, so taking global sections
  preserves injectivity on the integral scheme \(C\)).

  \paragraph{Lean signature scope.} The Lean definition
  \texttt{AlgebraicGeometry.Scheme.lineBundleAtClosedPoint.toFunctionField}
  takes the same inputs as \cref{def:lineBundleAtClosedPoint} (a curve
  \(C\), a point \(P\), and the closedness witness
  \texttt{IsClosed (\{P\} : Set C.left)}) and produces, for each global
  section \(s\) of \(\mathcal O_C(P)\) at twist \(0\) in the project's
  \(\Module \bar k\)-flavoured global-sections pipeline, the corresponding
  element \(\iota(s) \in K(C)\) of the project's function-field carrier
  \texttt{C.left.functionField}. The declaration is \texttt{noncomputable}
  and its body is gated on the body of
  \cref{def:lineBundleAtClosedPoint}.
\end{definition}

\section{Global sections as a Riemann--Roch space}

\begin{lemma}



\leanok
[Global sections of \(\mathcal O_C(P)\) as rational functions with a
controlled pole]
  \label{lem:lineBundleAtClosedPoint_globalSections_iff}
  \lean{AlgebraicGeometry.Scheme.lineBundleAtClosedPoint.globalSections_iff}
  
  % SOURCE: [Hartshorne], II.7, Proposition~7.7, p.~157 (effective
  % divisors linearly equivalent to \(D_0\) are the divisors of zeros of
  % global sections of \(\mathscr L(D_0)\), with the explicit
  % identification of those global sections with rational functions
  % \(f\) satisfying \((f) \geq -D_0\)) (read from
  % references/hartshorne-algebraic-geometry.pdf, PDF page 174).
  % SOURCE QUOTE: "Proposition 7.7. \emph{Let \(X\) be a nonsingular
  % projective variety over the algebraically closed field \(k\). Let
  % \(D_0\) be a divisor on \(X\) and let \(\mathscr L = \mathscr L(D_0)\)
  % be the corresponding invertible sheaf. Then:
  % (a) for each nonzero \(s \in \Gamma(X, \mathscr L)\), the divisor of
  % zeros \((s)_0\) is an effective divisor linearly equivalent to \(D_0\);
  % (b) every effective divisor linearly equivalent to \(D_0\) is \((s)_0\)
  % for some \(s \in \Gamma(X, \mathscr L)\); and
  % (c) two sections \(s, s' \in \Gamma(X, \mathscr L)\) have the same
  % divisor of zeros if and only if there is a \(\lambda \in k^*\) such
  % that \(s' = \lambda s\).}"
  % SOURCE QUOTE: "(b) If \(D \geq 0\) and \(D = D_0 + (f)\), then
  % \((f) \geq -D_0\). Thus \(f\) gives a global section of
  % \(\mathscr L(D_0)\) whose divisor of zeros is \(D\)."
  \textit{Source: Hartshorne, II.7, Proposition~7.7 and its
  proof of part (b), p.~157.}
  Let \(C\) be a smooth proper curve over \(\bar k\), let \(P \in C\) be a
  closed point, and let \(\mathcal O_C(P)\) be the line bundle of
  \cref{def:lineBundleAtClosedPoint}. The \(\bar k\)-vector space of
  global sections of \(\mathcal O_C(P)\) is naturally identified with the
  Riemann--Roch space
  \[
    L([P]) \;:=\;
    \bigl\{f \in K(C)^{\times} \;\big|\; \div(f) + [P] \geq 0
    \text{ in } \mathrm{Div}(C)\bigr\}
    \,\cup\, \{0\}.
  \]
  Equivalently, the order conditions
  \[
    \ord_Q(f) \;\geq\; 0 \quad \text{for every closed point } Q \in C
    \text{ with } Q \neq P,
    \qquad
    \ord_P(f) \;\geq\; -1
  \]
  on \(f \in K(C)^{\times}\) are equivalent to the existence of a global
  section \(s \in H^0(C, \mathcal O_C(P))\) whose image \(\iota(s)\) under
  the canonical inclusion $\iota : \mathcal O_C(P) \hookrightarrow
  \mathcal K_C \cong K(C)$ equals $f$ (the iff binds $s$ to $f$
  explicitly — vacuous-in-\(f\) readings such as ``some nonzero global
  section exists'' are excluded by construction).

  Under this identification, the inclusion
  \(\bar k \hookrightarrow H^0(C, \mathcal O_C(P))\) of constant functions
  corresponds to the natural map $\mathcal O_C \hookrightarrow
  \mathcal O_C(P)$ exhibiting $\mathcal O_C(P)$ as the twist of the
  structure sheaf along the prime divisor \([P]\).
\end{lemma}

\begin{proof}
  
  
  
  
  By construction, \(\mathcal O_C(P)\) is the sub-\(\mathcal O_C\)-module of
  the sheaf of total quotient rings \(\mathcal K_C\) generated locally by
  \(f_P^{-1}\) near \(P\) and by \(1\) on \(C \setminus \{P\}\), where \(f_P\) is
  a uniformizer at \(P\) (\cref{def:lineBundleAtClosedPoint}). Since \(C\)
  is integral, \(\mathcal K_C\) is the constant sheaf at the function field
  \(K(C)\), and a global section of any sub-\(\mathcal O_C\)-module
  \(\mathcal F \subset \mathcal K_C\) is precisely an element of \(K(C)\)
  whose germ at every point lies in \(\mathcal F\) at that point.
  Therefore \(f \in K(C)\) is a global section of \(\mathcal O_C(P)\) if
  and only if:

  \begin{itemize}
    \item For every closed point \(Q \in C\) with \(Q \neq P\), the germ of
      \(f\) at \(Q\) lies in \(\mathcal O_C(P)_Q = \mathcal O_{C, Q}\)
      (the structure sheaf, since \(\mathcal O_C(P)\) agrees with
      \(\mathcal O_C\) on \(C \setminus \{P\}\)); equivalently,
      \(\ord_Q(f) \geq 0\).
    \item At the closed point \(P\), the germ of \(f\) at \(P\) lies in
      \(\mathcal O_C(P)_P = f_P^{-1} \cdot \mathcal O_{C, P}\);
      equivalently, \(f_P \cdot f \in \mathcal O_{C, P}\), i.e.\
      \(\ord_P(f_P \cdot f) \geq 0\), i.e.\ \(\ord_P(f) \geq -1\).
  \end{itemize}
  In divisorial language (\cref{def:principal_divisor},
  \cref{def:divisor_closed_point}), the two conditions combine to
  \(\div(f) + [P] \geq 0\), recovering the Riemann--Roch space description
  \(L([P])\). The zero element of \(K(C)\) is included by convention as the
  zero global section.

  The identification with Hartshorne's Proposition~7.7(b) is direct:
  taking \(D_0 = [P]\) effective, Hartshorne's argument shows that a
  rational function \(f\) with \((f) \geq -D_0\) defines a global section
  of \(\mathscr L(D_0)\). The condition \((f) \geq -D_0 = -[P]\) is exactly
  \(\div(f) + [P] \geq 0\). The compatibility with the constant inclusion
  follows by taking \(f = 1 \in K(C)^{\times}\): then \(\div(1) = 0\), so
  \(\div(1) + [P] = [P] \geq 0\), and the corresponding section of
  \(\mathcal O_C(P)\) is the image of $1 \in H^0(C, \mathcal O_C)
  \cong \bar k$ under the structural inclusion $\mathcal O_C
  \hookrightarrow \mathcal O_C(P)$.
\end{proof}

\begin{lemma}\leanok
[Forward direction of \texttt{globalSections\_iff}: building a
section from the order conditions]
  \label{lem:lineBundleAtClosedPoint_globalSections_iff_mp}
  \lean{AlgebraicGeometry.Scheme.lineBundleAtClosedPoint.globalSections_iff_mp}
  \uses{lem:lineBundleAtClosedPoint_globalSections_iff,
        def:lineBundleAtClosedPoint,
        def:lineBundleAtClosedPoint_toFunctionField,
        def:lineBundleAtClosedPoint_carrierSubmoduleSheaf}
  Let \(C\) be a smooth proper geometrically irreducible integral curve over
  \(\bar k\), let \(P \in C\) be a closed point of coheight \(1\), and let
  \(f \in K(C)\) be nonzero satisfying the order conditions \(\ord_Q(f) \ge 0\)
  for every prime divisor \(Q \ne P\) and \(\ord_P(f) \ge -1\). Then there
  exists a global section \(s \in H^0(C, \mathcal O_C(P))\) whose image under
  the canonical inclusion \(\iota\) into \(K(C)\) equals \(f\). This is the
  forward (\texttt{mp}) direction of
  \cref{lem:lineBundleAtClosedPoint_globalSections_iff}.
\end{lemma}

\begin{proof}
  \uses{def:lineBundleAtClosedPoint_carrierSubmoduleSheaf,
        def:lineBundleAtClosedPoint_toFunctionField}
  Proved directly in Lean. The order conditions place \(f\) in
  \(\mathrm{carrierSubmoduleSheaf}_P(\top)\); the resulting witness spans a
  \(\bar k\)-linear map \(\bar k \to \mathrm{carrierSubmoduleSheaf}_P(\top)\),
  which transports along the constant-sheaf adjunction and the
  global-sections linear equivalence to a section \(s\) with
  \(\iota(s) = f\).
\end{proof}

\section{Cohomological vanishing in genus zero}

\begin{lemma}



\leanok
[Vanishing of \(H^1(C, \mathcal O_C(P))\) on a genus-\(0\) curve]
  \label{lem:H1_vanishing_lineBundleAtClosedPoint_genusZero}
  \lean{AlgebraicGeometry.Scheme.lineBundleAtClosedPoint.h1_vanishing_genusZero}
  \uses{def:lineBundleAtClosedPoint, def:Scheme_HModule, def:genus,
        thm:H1_vanishing_flasque, lem:sheafOf_ses_single_add}

  % SOURCE: [Hartshorne], IV.1, the proof of Theorem~1.3 (the inductive
  % step ``Next, let \(D\) be any divisor, and let \(P\) be any point.''):
  % the short exact sequence $0 \to \mathscr L(-P) \to \mathscr O_X \to
  % k(P) \to 0\( and its tensor with \)\mathscr L(D + P)$ giving
  % \(0 \to \mathscr L(D) \to \mathscr L(D + P) \to k(P) \to 0\)
  % (read from references/hartshorne-algebraic-geometry.pdf, PDF page 313).
  % SOURCE QUOTE: "We consider \(P\) as a closed subscheme of \(X\). Its
  % structure sheaf is a skyscraper sheaf \(k\) sitting at the point \(P\),
  % which we denote by \(k(P)\), and its ideal sheaf is \(\mathscr L(-P)\)
  % by (II, 6.18). Therefore we have an exact sequence
  % \(0 \to \mathscr L(-P) \to \mathscr O_X \to k(P) \to 0\).
  % Tensoring with \(\mathscr L(D + P)\) we get
  % \(0 \to \mathscr L(D) \to \mathscr L(D + P) \to k(P) \to 0\),
  % (since \(\mathscr L(P)\) is locally free of rank 1, tensoring by it
  % does not affect the sheaf \(k(P)\))."
  \textit{Source: Hartshorne, IV.1, p.~296, the inductive step of the
  proof of Theorem~1.3.}
  Let \(C\) be a smooth proper geometrically irreducible curve over
  \(\bar k\) with \(g(C) = 0\), and let \(P \in C\) be a closed point. Then
  \[
    H^1(C, \mathcal O_C(P)) \;=\; 0,
  \]
  i.e.\ \(\dim_{\bar k} H^1(C, \mathcal O_C(P)) = 0\).
\end{lemma}

\begin{proof}
  
  Specialise Hartshorne's inductive step at \(D = 0\). The standard short
  exact sequence
  \[
    0 \to \mathcal O_C(-[P]) \to \mathcal O_C \to k(P) \to 0
  \]
  (Hartshorne~II.6.18: the ideal sheaf of the locally principal closed
  subscheme \(P\) is \(\mathcal O_C(-[P])\); the quotient is the structure
  sheaf of \(P\), which is the skyscraper sheaf \(k(P) \cong \bar k\)
  supported at \(P\) since \(P\) has residue field \(\bar k\) by closedness
  over the algebraically closed base \(\bar k\)). Tensoring with the
  locally free rank-\(1\) sheaf \(\mathcal O_C([P])\)
  (\cref{def:lineBundleAtClosedPoint}) preserves exactness and leaves
  the skyscraper sheaf \(k(P)\) unchanged, giving
  \[
    0 \to \mathcal O_C \to \mathcal O_C(P) \to k(P) \to 0
  \]
  in \(\mathbf{Coh}(C)\). Apply the long exact sequence of sheaf
  cohomology (\cref{def:Scheme_HModule}, the project's $\Module \bar
  k$-flavoured cohomology pipeline). Using that $\dim C = 1$ kills
  \(H^i\) for \(i \geq 2\) (Grothendieck vanishing on a \(1\)-dimensional
  scheme), the long exact sequence reads
  \begin{multline*}
    0 \to H^0(C, \mathcal O_C) \to H^0(C, \mathcal O_C(P)) \to H^0(C, k(P)) \\
    \to H^1(C, \mathcal O_C) \to H^1(C, \mathcal O_C(P)) \to H^1(C, k(P)) \to 0.
  \end{multline*}
  Two ingredients close the proof:

  \emph{(a) \(H^1(C, \mathcal O_C) = 0\).} This is the genus-\(0\)
  hypothesis: by \cref{def:genus},
  \[
    g(C) \;=\; \dim_{\bar k} H^1(C, \mathcal O_C),
  \]
  so \(g(C) = 0\) is exactly \(\dim_{\bar k} H^1(C, \mathcal O_C) = 0\),
  i.e.\ \(H^1(C, \mathcal O_C) = 0\) as a \(\bar k\)-vector space.

  \emph{(b) \(H^1(C, k(P)) = 0\).} The skyscraper sheaf \(k(P)\) is
  supported at a single closed point, hence is flasque: every section
  on an open \(U \subset C\) either restricts to \(0\) (if \(P \notin U\)) or
  takes values in the stalk \(k(P)_P \cong \bar k\) (if \(P \in U\)), and
  restriction maps between opens of either type are surjective. A
  flasque sheaf on any topological space has vanishing higher
  cohomology (Hartshorne~III.2.5); in particular
  \(H^i(C, k(P)) = 0\) for \(i \geq 1\), and \emph{a fortiori}
  \(H^1(C, k(P)) = 0\).

  Substituting (a) and (b) into the long exact sequence collapses the
  segment ending at \(H^1(C, \mathcal O_C(P))\) to
  \[
    0 \to H^1(C, \mathcal O_C(P)) \to 0,
  \]
  whence \(H^1(C, \mathcal O_C(P)) = 0\), as required.

  \emph{Lean reference note on the long exact sequence.} Mathlib
  (snapshot \texttt{b80f227}) ships the long exact sequence of
  \(\Ext_{\mathcal O_C}^*(\mathbf 1, -)\) at the abelian-group level
  (\texttt{CategoryTheory.Abelian.Ext.covariantSequence\_exact}); the
  project's \(\Module \bar k\)-flavoured cohomology pipeline
  (\cref{def:Scheme_HModule}) inherits the long exact sequence by
  forget-functor naturality. The skyscraper-cohomology vanishing
  \(H^i(C, k(P)) = 0\) for \(i \geq 1\) is closed by either the flasque
  argument (Hartshorne~III.2.5; Mathlib has the abelian-sheaf form
  via \texttt{CategoryTheory.Sheaf.Flasque}) or, more directly for the
  curve case, by the affine-cover Čech-cohomology argument on
  \(C \setminus \{P\}\) and a small affine open around \(P\) (the
  skyscraper has zero sections on \(C \setminus \{P\}\), so the Čech
  complex degenerates).
\end{proof}

\section{The dimension formula in genus zero}

% NOTE (iter-194 reviewer): The Lean body is now structurally
% sorry-free (closed transitively through two named substrate
% helpers). The substrate helpers `h0_sub_h1_lineBundleAtClosedPoint_eq_two`
% (the Int-valued $H^0 - H^1 = 2$ \(\chi\)-bridge; gated on
% `OcOfD.lean` STRUCTURALLY BLOCKED + `RRFormula.lean` $\chi$-additivity)
% and `h1_vanishing_genusZero` (gated on Lane H \texttt{H1Vanishing.lean}
% substrate) carry the remaining typed sorries.  Future consumers
% reading the body should consult those.
\begin{theorem}



\leanok
[\(\dim_{\bar k} H^0(C, \mathcal O_C(P)) = 2\) on a genus-\(0\) curve]
  \label{thm:lineBundleAtClosedPoint_dim_eq_two_of_genusZero}
  \lean{AlgebraicGeometry.Scheme.lineBundleAtClosedPoint.dim_eq_two_of_genusZero}
  \uses{thm:euler_char_eq_deg_plus_one_minus_genus, def:eulerChar_curve,
        def:divisor_closed_point, def:divisor_degree, def:genus,
        def:lineBundleAtClosedPoint,
        thm:lineBundleAtClosedPoint_h0_sub_h1_eq_two,
        lem:H1_vanishing_lineBundleAtClosedPoint_genusZero}

  % SOURCE: [Hartshorne], IV.1, the Riemann--Roch chain inside
  % Example~1.3.5 applied with \(D = [P]\): $\chi(\mathscr L([P]))
  % = \deg([P]) + 1 - g = 1 + 1 - 0 = 2$, combined with the vanishing
  % \(H^1(\mathscr L([P])) = 0\) (read from
  % references/hartshorne-algebraic-geometry.pdf, PDF pages 312, 314).
  % SOURCE QUOTE: "Theorem 1.3 (Riemann--Roch). \emph{Let \(D\) be a
  % divisor on a curve \(X\) of genus \(g\). Then}
  % \(l(D) - l(K - D) = \deg D + 1 - g\)."
  % SOURCE QUOTE: "Thus we have to show that for any \(D\),
  % \(\chi(\mathscr L(D)) = \deg D + 1 - g\),
  % where for any coherent sheaf \(\mathscr F\) on \(X\), \(\chi(\mathscr F)\)
  % is the Euler characteristic
  % \(\chi(\mathscr F) = \dim H^0(X, \mathscr F) - \dim H^1(X, \mathscr F)\)."
  \textit{Source: Hartshorne, IV.1, Theorem~1.3, p.~295 (the
  \(\chi\)-identity inside the proof) and Example~1.3.5, p.~297 (the
  genus-\(0\) specialisation chain).}
  Let \(C\) be a smooth proper geometrically irreducible curve over
  \(\bar k\) with \(g(C) = 0\), and let \(P \in C\) be a closed point. Then
  \[
    \dim_{\bar k} H^0(C, \mathcal O_C(P)) \;=\; 2.
  \]
\end{theorem}

\begin{proof}
  
  
  
  Specialise the Euler-characteristic identity
  \cref{thm:euler_char_eq_deg_plus_one_minus_genus} of \texttt{RR.2} to
  the divisor \(D = [P]\) of \cref{def:divisor_closed_point}. The
  degree of \([P]\) is \(\deg([P]) = 1\) (every closed point contributes
  \(1\) to the bare-sum degree of \cref{def:divisor_degree}), and the
  genus hypothesis \(g(C) = 0\) (\cref{def:genus}) gives
  \[
    \chi(\mathcal O_C(P)) \;=\; \deg([P]) + 1 - g(C)
    \;=\; 1 + 1 - 0 \;=\; 2.
  \]
  Unfolding \(\chi\) via \cref{def:eulerChar_curve} as a two-term Euler
  characteristic on a curve, this rewrites as
  \[
    \dim_{\bar k} H^0(C, \mathcal O_C(P))
    \,-\, \dim_{\bar k} H^1(C, \mathcal O_C(P))
    \;=\; 2.
  \]
  By
  \cref{lem:H1_vanishing_lineBundleAtClosedPoint_genusZero}, the second
  term vanishes:
  \(\dim_{\bar k} H^1(C, \mathcal O_C(P)) = 0\). Substituting yields
  \(\dim_{\bar k} H^0(C, \mathcal O_C(P)) = 2\), as required.

  \emph{Alternative direct route via the long exact sequence.} The
  same long exact sequence used in the proof of
  \cref{lem:H1_vanishing_lineBundleAtClosedPoint_genusZero} computes
  the dimension of \(H^0\) in one stroke without invoking
  \texttt{RR.2}: from
  \[
    0 \to H^0(C, \mathcal O_C) \to H^0(C, \mathcal O_C(P))
    \to H^0(C, k(P)) \to H^1(C, \mathcal O_C) \to \cdots
  \]
  together with \(\dim_{\bar k} H^0(C, \mathcal O_C) = 1\)
  (Hartshorne~I.3.4: \(H^0\) of the structure sheaf on a proper
  geometrically irreducible variety over \(\bar k\) is \(\bar k\)),
  \(\dim_{\bar k} H^0(C, k(P)) = 1\) (skyscraper at a \(\bar k\)-rational
  point), and \(H^1(C, \mathcal O_C) = 0\) (genus-\(0\)), the connecting
  map \(H^0(C, k(P)) \to H^1(C, \mathcal O_C)\) has zero codomain and the
  segment yields the short exact sequence
  \[
    0 \to \bar k \to H^0(C, \mathcal O_C(P)) \to \bar k \to 0
  \]
  of \(\bar k\)-vector spaces, whence $\dim_{\bar k}
  H^0(C, \mathcal O_C(P)) = 1 + 1 = 2$. This alternative route does
  not consume the Euler-characteristic identity of \texttt{RR.2} and
  could be substituted at the formalisation stage; the chapter uses
  \texttt{RR.2} as written above to keep the dependency graph in
  line with the iter-173 RR sub-build sequencing.
\end{proof}

\begin{theorem}\leanok
[The Euler-characteristic bridge \(H^0 - H^1 = 2\) at
\(D = [P]\)]
  \label{thm:lineBundleAtClosedPoint_h0_sub_h1_eq_two}
  \lean{AlgebraicGeometry.Scheme.lineBundleAtClosedPoint.h0_sub_h1_lineBundleAtClosedPoint_eq_two}
  \uses{def:lineBundleAtClosedPoint,
        thm:euler_char_eq_deg_plus_one_minus_genus,
        def:divisor_degree, def:genus}
  Let \(C\) be a smooth proper geometrically irreducible curve over \(\bar k\)
  with \(g(C) = 0\), and let \(P \in C\) be a closed point of coheight \(1\).
  Then, as integers,
  \[
    \dim_{\bar k} H^0(C, \mathcal O_C(P))
    \,-\, \dim_{\bar k} H^1(C, \mathcal O_C(P)) \;=\; 2.
  \]
  This is the specialisation of the Euler-characteristic identity
  \(\chi(\mathcal O_C(D)) = \deg(D) + 1 - g(C)\) to the single-point divisor
  \(D = [P]\), with \(\deg([P]) = 1\) and \(g(C) = 0\); it isolates the
  \(\chi\)-arithmetic consumed by
  \cref{thm:lineBundleAtClosedPoint_dim_eq_two_of_genusZero}.
\end{theorem}

\begin{proof}
  Proved directly in Lean, as a sub-step of
  \cref{thm:lineBundleAtClosedPoint_dim_eq_two_of_genusZero} via the
  Riemann--Roch \(\chi\)-identity of \texttt{RR.2} evaluated at the
  single-point divisor \([P]\).
\end{proof}

\subsection{Iter-195 substrate substeps for the
  non-constant-rational corollary}
\label{subsec:lineBundleAtClosedPoint_substrate_substeps}

The iter-195 substrate prover reduced the body of
\cref{cor:lineBundleAtClosedPoint_exists_nonconstant_genusZero} to a
private helper \texttt{exists\_nonconstant\_rational\_from\_dim\_eq\_two}
(file \texttt{AlgebraicJacobian/RiemannRoch/OCofP.lean}, line~1323)
which carries the substantive linear-algebra plus principal-divisor
content of the corollary, isolated from the genus hypothesis. The
helper's body, after the iter-195 push, executes six structural setup
steps culminating in the choice of a non-constant candidate section
\(s\in H^0(C, \mathcal O_C(P))\setminus \bar k\cdot s_1\) (with
\(s_1\) the distinguished constant section satisfying
\(\iota(s_1)=1\)) and the definition
\[
  f \;:=\; \iota(s) \;=\;
  \texttt{lineBundleAtClosedPoint.toFunctionField}\,P\,h_P\,h_{P\mathrm{coh}}\,s
  \;\in\; K(C).
\]
After this setup, three named sub-claims (a), (b), (c) remain to close
the helper. The iter-196 plan splits each into a project-bespoke
substrate lemma carrying its own \texttt{\textbackslash lean}-pin and proof recipe. The
sub-claims are Archon-original decompositions of the helper's
in-Lean documentation block at lines~1421--1440 of
\texttt{OCofP.lean}; the underlying mathematics is the standard
content of Hartshorne~IV.1 Exercise~1.1 (already cited at the
\(\cref{cor:lineBundleAtClosedPoint_exists_nonconstant_genusZero}\)
level) augmented by the Stacks Project tag \texttt{02P0} input cited
inline in sub-claim~(c).

\begin{theorem}\leanok
[Existence of a non-constant rational function from
\(\dim H^0 = 2\)]
  \label{thm:lineBundleAtClosedPoint_exists_nonconstant_rational_from_dim_eq_two}
  \lean{AlgebraicGeometry.Scheme.lineBundleAtClosedPoint.exists_nonconstant_rational_from_dim_eq_two}
  \uses{def:lineBundleAtClosedPoint,
        lem:lineBundleAtClosedPoint_globalSections_iff_mp,
        lem:lineBundleAtClosedPoint_globalSections_iff_mpr,
        lem:lineBundleAtClosedPoint_toFunctionField_injective,
        lem:lineBundleAtClosedPoint_functionField_const_of_complete_curve_of_orderZero,
        def:principal_divisor}
  Let \(C\) be a smooth proper geometrically irreducible integral curve over
  \(\bar k\), let \(P \in C\) be a closed point of coheight \(1\), and suppose
  \(\dim_{\bar k} H^0(C, \mathcal O_C(P)) = 2\). Then there exists a nonzero
  rational function \(f \in K(C)\) satisfying \(\ord_Q(f) \ge 0\) for every
  prime divisor \(Q \ne P\), \(\ord_P(f) \ge -1\), and with nonzero principal
  divisor \(\div(f) \ne 0\). This is the genus-free substrate carrying the
  linear-algebra and principal-divisor content of
  \cref{cor:lineBundleAtClosedPoint_exists_nonconstant_genusZero}, isolated
  from the genus hypothesis (only the dimension count is consumed).
\end{theorem}

\begin{proof}
  Proved directly in Lean. The two-dimensionality of
  \(H^0(C, \mathcal O_C(P))\) together with the one-dimensional constant
  subspace \(\bar k \cdot s_1\) (with \(s_1\) the distinguished constant
  section, \(\iota(s_1) = 1\)) yields a section \(s \notin \bar k \cdot s_1\).
  Setting \(f := \iota(s)\), the inclusion is injective by
  \cref{lem:lineBundleAtClosedPoint_toFunctionField_injective} (so
  \(f \ne 0\)), the order conditions on \(f\) follow from
  \cref{lem:lineBundleAtClosedPoint_globalSections_iff_mpr}, and the
  non-vanishing of \(\div(f)\) follows by contraposition from
  \cref{lem:lineBundleAtClosedPoint_functionField_const_of_complete_curve_of_orderZero}
  (were \(\div(f) = 0\), the function would be constant, forcing
  \(s \in \bar k \cdot s_1\)).
\end{proof}

\begin{lemma}


\leanok
[Sub-claim (a): \(\texttt{toFunctionField}\) is \(\bar k\)-linearly
injective, hence preserves non-zero-ness of sections]
  \label{lem:lineBundleAtClosedPoint_toFunctionField_injective}
  \lean{AlgebraicGeometry.Scheme.lineBundleAtClosedPoint.toFunctionField_injective}
  \uses{def:lineBundleAtClosedPoint_toFunctionField,
        def:lineBundleAtClosedPoint_carrierSubmoduleSheaf,
        def:lineBundleAtClosedPoint,
        lem:lineBundleAtClosedPoint_globalSections_iff,
        def:Scheme_HModule_zero_linearEquiv}
  Let \(C\) be a smooth proper geometrically irreducible integral curve
  over \(\bar k\) and let \(P\in C\) be a closed point with coheight
  witness \(h_{P\mathrm{coh}}\colon \mathrm{coheight}(P)=1\). Then the
  \(\bar k\)-linear map
  \[
    \iota \;:=\;
    \texttt{lineBundleAtClosedPoint.toFunctionField}\,P\,h_P\,h_{P\mathrm{coh}}
    \;\colon\;
    H^0(C, \mathcal O_C(P)) \;\longrightarrow\; K(C)
  \]
  of \cref{def:lineBundleAtClosedPoint_toFunctionField} is injective.
  In particular, for any global section \(s\in H^0(C, \mathcal O_C(P))\)
  with \(s\ne 0\), the rational function \(f:=\iota(s)\in K(C)\) is
  non-zero.
\end{lemma}

\begin{proof}
  \uses{lem:lineBundleAtClosedPoint_globalSections_iff,
        def:lineBundleAtClosedPoint_carrierSubmoduleSheaf}
  The map \(\iota\) factors as the composition of four \(\bar k\)-linear
  isomorphisms and one final injective \(\bar k\)-linear map; each
  factor's injectivity is delivered by existing Mathlib or
  project-substrate hooks. Concretely, on the integral scheme \(C\),
  \(\iota\) equals the composition:
  \begin{enumerate}
    \item \(\texttt{HModule\_zero\_linearEquiv}_{\bar k,\,\mathcal O_C(P)}\):
      the project's $\bar k$-linear isomorphism
      $H^0(C, \mathcal O_C(P)) \;=\; \texttt{Scheme.HModule}\,
      \bar k\,\mathcal O_C(P)\,0
        \;\xrightarrow{\sim}\;
        \Hom_{\Sh}(\mathbf 1, \mathcal O_C(P))$
      with \(\mathbf 1\) the constant sheaf of \(\bar k\) on \(C\). This
      bridge is the project-substrate
      \texttt{Cohomology/StructureSheafModuleK.lean} entry
      \texttt{HModule\_zero\_linearEquiv}.
    \item \texttt{sheafToPresheaf.map}: the fully faithful functor
      \(\Sheaf \to \Psh\) on sheaves of \(\bar k\)-modules sends the
      sheaf-hom of layer~(1) to the underlying presheaf-hom; full
      faithfulness is Mathlib's
      \texttt{CategoryTheory.Sheaf.fullyFaithfulSheafToPresheaf}, so
      this step is injective on hom-sets.
    \item \(\texttt{constantSheafAdj}_{\bar k,\,\mathcal O_C(P)}.\texttt{homEquiv}\):
      Mathlib's constant-sheaf-adjunction
      \(\texttt{constantSheaf} \dashv \Gamma\) supplies the
      \(\bar k\)-linear isomorphism
      $\Hom_{\Psh}(\mathbf 1_{\mathrm{psh}}, \mathcal O_C(P).\val)
       \;\xrightarrow{\sim}\;
       \Hom_{\ModuleCat\,\bar k}(\bar k, \mathcal O_C(P).\val(\top))$.
    \item Evaluation at the generator \(1\in\bar k\): the canonical
      \(\bar k\)-linear isomorphism
      $\Hom_{\ModuleCat\,\bar k}(\bar k, M)
       \;\xrightarrow{\sim}\;
       M, \qquad \varphi \mapsto \varphi(1),$
      with \(M:=\mathcal O_C(P).\val(\top)\) the global-section module
      of \cref{def:lineBundleAtClosedPoint_carrierSubmoduleSheaf}. The
      Mathlib name at the project's pinned commit may be
      \texttt{LinearMap.applyOneEquiv} or its
      \texttt{LinearEquiv.fromBaseRing} variant; the prover hand-rolls
      a 5-line \texttt{LinearEquiv} if the exact name has been
      renamed upstream.
    \item \texttt{Subtype.val}: the carrier sheaf
      \(\mathcal O_C(P).\val(\top) \;=\;
      \texttt{carrierSubmoduleSheaf}\,P\,h_{P\mathrm{coh}}(\op\,\top)\)
      is a \(\bar k\)-submodule of the function field \(K(C)\)
      (\cref{def:lineBundleAtClosedPoint_carrierSubmoduleSheaf}); the
      coercion \texttt{Subtype.val} to \(K(C)\) is injective.
  \end{enumerate}
  Each layer (1)--(4) is a \texttt{LinearEquiv}; layer (5) is
  \texttt{Function.Injective Subtype.val}. The composition is therefore
  an injective \(\bar k\)-linear map, and this composition is precisely
  the body unfolding of
  \cref{def:lineBundleAtClosedPoint_toFunctionField}: the iter-187
  refactor of \texttt{toFunctionField} is the formal record of this
  factorisation (compare the \emph{iter-187 closure} commentary at
  \cref{def:lineBundleAtClosedPoint_toFunctionField} and the
  \texttt{globalSections\_iff\_mpr} proof at line~1024 of
  \texttt{OCofP.lean}, where the same chain is unfolded in the
  reverse direction). The non-zero-transfer conclusion is immediate:
  if \(\iota(s)=0\) and \(s\ne 0\), then \(\iota\) has a non-trivial
  kernel, contradicting injectivity; equivalently, applying
  injectivity to \(\iota(s)=0=\iota(0)\) (using
  \texttt{htF\_zero}, already in scope at line~1364 of the helper
  body) gives \(s=0\), contradicting the iter-195 hypothesis
  \(s\notin \bar k\cdot s_1\) (which is strictly stronger than
  \(s\ne 0\)).

  \emph{Estimated formalisation footprint.} \(\sim 30\)--\(50\)
  Lean LOC: layer~(1) is one line (named bridge); layer~(2) is
  \(\sim 5\) LOC; layer~(3) is \(\sim 5\)--\(10\) LOC; layer~(4) is
  \(\sim 5\)--\(10\) LOC (hand-roll if necessary); layer~(5) is one
  line (\texttt{Subtype.val\_injective}); the composition and the
  final \(s\ne 0\Rightarrow\iota(s)\ne 0\) wrap each add 5--10 LOC.
\end{proof}

\begin{lemma}


\leanok
[Sub-claim (b): backward direction of \texttt{globalSections\_iff} —
order conditions from an existing section]
  \label{lem:lineBundleAtClosedPoint_globalSections_iff_mpr}
  \lean{AlgebraicGeometry.Scheme.lineBundleAtClosedPoint.globalSections_iff_mpr}
  \uses{lem:lineBundleAtClosedPoint_globalSections_iff,
        def:lineBundleAtClosedPoint_toFunctionField,
        def:lineBundleAtClosedPoint_carrierSubmoduleSheaf}
  Let \(C, P, h_P, h_{P\mathrm{coh}}\) be as in
  \cref{lem:lineBundleAtClosedPoint_toFunctionField_injective}. Let
  \(f \in K(C)\) be nonzero, and suppose there exists a global section
  \(s \in H^0(C, \mathcal O_C(P))\) with \(\iota(s) = f\), where
  \(\iota:=
  \texttt{lineBundleAtClosedPoint.toFunctionField}\,P\,h_P\,h_{P\mathrm{coh}}\).
  Then \(f\) satisfies the order conditions:
  \[
    \ord_Q(f) \;\ge\; 0
    \qquad \text{for every prime divisor } Q\in C \text{ with } Q\ne P,
    \qquad
    \ord_P(f) \;\ge\; -1.
  \]
  This is the backward direction (\texttt{mpr}) of the equivalence
  \cref{lem:lineBundleAtClosedPoint_globalSections_iff}; it is pinned as
  a named substrate helper so that
  \cref{cor:lineBundleAtClosedPoint_exists_nonconstant_genusZero}'s Lean
  body can invoke it directly with a witness
  \(\langle s, \mathtt{rfl}\rangle\) built from the iter-195 setup.
\end{lemma}

\begin{proof}
  \uses{def:lineBundleAtClosedPoint_carrierSubmoduleSheaf,
        def:lineBundleAtClosedPoint_toFunctionField}
  Unfold the witness section \(s\) along the linear-equivalence chain
  defining \(\iota = \texttt{toFunctionField}\) (the five-layer
  composition recorded in
  \cref{lem:lineBundleAtClosedPoint_toFunctionField_injective}). The
  underlying carrier element at the top open is a pair
  \((\widetilde f,\,p) \in \mathrm{carrierSubmoduleSheaf}_P(\op\top)\)
  with \(\widetilde f = f\) (from the hypothesis \(\iota(s) = f\)) and
  \(p\) a membership witness. By
  \cref{def:lineBundleAtClosedPoint_carrierSubmoduleSheaf}, membership
  in \(\mathrm{carrierSubmoduleSheaf}_P(\op\top)\) unfolds as the
  conjunction of (i) order \(\ge 0\) at every prime divisor \(Q\) of
  \(\top\) other than \(P\), and (ii) order \(\ge -1\) at \(P\)
  (intersected with the trivial \texttt{trivAtBot} factor, which is
  \(\top\) on the top open). Projecting \(p\) onto its first and
  second components returns exactly the pair of order bounds claimed
  (the auxiliary \(Q.\mathrm{point}\in \mathrm{Set.univ}\) premises
  are discharged by \(\mathrm{Set.mem\_univ}\)). The nonzero hypothesis
  on \(f\) is carried for downstream typing but is not consumed in the
  membership-projection step.
\end{proof}

\begin{lemma}


\leanok
[Sub-claim (c) substrate: a rational function with zero order at every
prime divisor on a complete curve is constant
(Hartshorne~I.3.4 / Stacks~\texttt{02P0})]
  \label{lem:lineBundleAtClosedPoint_functionField_const_of_complete_curve_of_orderZero}
  \lean{AlgebraicGeometry.Scheme.lineBundleAtClosedPoint.functionField_const_of_complete_curve_of_orderZero}
  \uses{def:lineBundleAtClosedPoint, lem:lineBundleAtClosedPoint_algebraMap_bijective_of_finite_isDomain_isAlgClosed, lem:lineBundleAtClosedPoint_functionField_localUnit_of_orderZero_at_primeDivisor}
  Let \(C\) be a smooth proper geometrically irreducible integral curve
  over the algebraically closed field \(\bar k\). Let \(f \in K(C)\) be
  a nonzero rational function such that
  \(\ord_Q(f) = 0\) for every prime divisor \(Q \in C\). Then \(f\)
  lies in the image of the constant inclusion
  \(\bar k \hookrightarrow K(C)\): there exists \(c \in \bar k\) with
  \[
    f \;=\; \texttt{algebraMap}\,\bar k\,K(C)\,c.
  \]
\end{lemma}

% NOTE (iter-197 rename — substrate vs.\ contrapositive). The original
% iter-196 phrasing of sub-claim~(c) was the contrapositive form
% ``principal divisor of a non-constant rational function on a complete
% curve is non-zero'':
%   \mathtt{Scheme.WeilDivisor.principal}\,f\,h_f \;\ne\; 0
% whenever \(s \notin \bar k\cdot s_1\) and \(f = \iota(s)\). The two
% forms are mathematically equivalent: \(\mathrm{principal}\,f =
% \sum_Q \ord_Q(f)\cdot[Q]\), so \(\mathrm{principal}\,f = 0\) iff
% \(\ord_Q(f) = 0\) for every prime divisor \(Q\), and the substrate
% above then forces \(f \in \bar k\). The iter-196 prover chose to
% extract the Hartshorne~I.3.4 / Stacks~\texttt{02P0} form as the
% standalone named substrate — it cleanly isolates the substantive
% Mathlib gap (``\(\Gamma(C, \mathcal O_C) = \bar k\) on a complete
% geometrically irreducible curve over \(\bar k\)'') — and inlined the
% contrapositive ``\(s \notin \bar k\cdot s_1 \Rightarrow
% \mathrm{principal}\,f \ne 0\)'' wrapper inside
% \texttt{exists\_nonconstant\_rational\_from\_dim\_eq\_two}
% (\texttt{AlgebraicJacobian/RiemannRoch/OCofP.lean} lines~1596--1635).
% The contrapositive runs as follows: assume \(\mathrm{principal}\,f
% = 0\); apply the substrate to obtain \(f = c \in \bar k\); use
% \(\bar k\)-linearity of \(\iota\) plus \(\iota(s_1) = 1\) and
% injectivity of \(\iota\)
% (\cref{lem:lineBundleAtClosedPoint_toFunctionField_injective}) to
% conclude \(s = c\cdot s_1 \in \bar k\cdot s_1\), contradicting
% \(s \notin \bar k\cdot s_1\).

\begin{proof}
  \uses{def:lineBundleAtClosedPoint, lem:lineBundleAtClosedPoint_functionField_localUnit_of_orderZero_at_primeDivisor, lem:lineBundleAtClosedPoint_algebraMap_bijective_of_finite_isDomain_isAlgClosed}
  This is the classical Hartshorne~I.3.4 statement
  ``\(\Gamma(C, \mathcal O_C) = \bar k\) on a complete geometrically
  irreducible reduced variety over the algebraically closed field
  \(\bar k\)'', combined with the
  ``no-zeros-no-poles \(\Rightarrow\) regular'' content of
  Stacks~\texttt{02P0}, in two ingredients:

  \emph{Step (i) — algebraic Hartogs.} On the normal Noetherian scheme
  \(C\), a rational function regular at every codimension-\(1\) point
  extends uniquely to a global section of the structure sheaf
  (Stacks tag~\texttt{0BCK}: $\Gamma(C, \mathcal O_C) =
  \bigcap_{Q\in C^{(1)}} \mathcal O_{C, Q}$). On the smooth curve
  \(C\) every prime divisor \(Q\) is a codimension-\(1\) point and
  \(\ord_Q\) is the standard discrete valuation on
  \(\mathcal O_{C, Q}\); the hypothesis \(\ord_Q(f) = 0\) for every
  \(Q\) places \(f\) in the intersection
  \(\bigcap_Q \mathcal O_{C, Q} = \Gamma(C, \mathcal O_C)\).

  \emph{Step (ii) — global sections are constants.} On a proper
  geometrically irreducible reduced variety over the algebraically
  closed field \(\bar k\), the global-sections ring of the structure
  sheaf coincides with the base field,
  \(\Gamma(C, \mathcal O_C) = \bar k\) (Hartshorne~I.3.4; or
  cohomologically: \(\Gamma\) is finite-dimensional over \(\bar k\)
  by Hartshorne~III.5.2, hence integral over \(\bar k\), hence
  algebraic, hence \(\bar k\) by algebraic closure). Combining
  step~(i) with step~(ii) places \(f\) in the image of the constant
  inclusion \(\bar k \hookrightarrow K(C)\), as required.

  \emph{Mathlib status.} At snapshot \texttt{b80f227} both ingredients
  are project-bespoke Mathlib gaps: the Lean substrate helper carries
  a single typed sorry pending upstream development; the named-helper
  extraction makes the upstream dependency explicit for future
  iterations.
\end{proof}

\begin{lemma}\leanok
[Local DVR lift from vanishing order]
  \label{lem:lineBundleAtClosedPoint_localLift_of_log_ordFrac_eq_zero}
  \lean{AlgebraicGeometry.Scheme.lineBundleAtClosedPoint.localLift_of_log_ordFrac_eq_zero}
  Let \(R\) be a discrete valuation ring with field of fractions \(K\). A
  nonzero element \(x \in K\) whose valuation is \(0\) (equivalently,
  \(\log\) of its order is \(0\)) is the image of a unit of \(R\): there
  exists a unit \(r \in R^{\times}\) with \(x = r\) in \(K\).
\end{lemma}

\begin{proof}
  Proved directly in Lean: vanishing of \(\log\) of the order forces the
  order to equal \(1\), placing \(x\) in the multiplicative kernel of the
  order homomorphism, which for a discrete valuation ring is exactly the
  image of the unit submonoid.
\end{proof}

\begin{lemma}\leanok
[A finite integral-domain algebra over an algebraically closed
field is the field]
  \label{lem:lineBundleAtClosedPoint_algebraMap_bijective_of_finite_isDomain_isAlgClosed}
  \lean{AlgebraicGeometry.Scheme.lineBundleAtClosedPoint.algebraMap_bijective_of_finite_isDomain_isAlgClosed}
  Let \(k\) be an algebraically closed field and \(R\) an integral domain that
  is a finite-dimensional \(k\)-algebra. Then the structural algebra map
  \(k \to R\) is bijective. This is the abstract algebraic kernel of the
  identity \(\Gamma(C, \mathcal O_C) = \bar k\) on a complete curve.
\end{lemma}

\begin{proof}
  Proved directly in Lean, as a re-export of Mathlib's
  \texttt{IsAlgClosed.algebraMap\_bijective\_of\_isIntegral} against the
  finite-implies-integral instance.
\end{proof}

\begin{lemma}\leanok
[Stalk-unit lift from vanishing order at a prime divisor]
  \label{lem:lineBundleAtClosedPoint_functionField_localUnit_of_orderZero_at_primeDivisor}
  \lean{AlgebraicGeometry.Scheme.lineBundleAtClosedPoint.functionField_localUnit_of_orderZero_at_primeDivisor}
  \uses{lem:lineBundleAtClosedPoint_localLift_of_log_ordFrac_eq_zero}
  Let \(X\) be an integral, locally Noetherian, regular-in-codimension-one
  scheme, let \(Q\) be a prime divisor of \(X\), and let \(f \in K(X)\) be
  nonzero with \(\ord_Q(f) = 0\). Then \(f\) is the image of a unit of the
  stalk \(\mathcal O_{X, Q}\) under the canonical embedding
  \(\mathcal O_{X, Q} \hookrightarrow K(X)\).
\end{lemma}

\begin{proof}
  \uses{lem:lineBundleAtClosedPoint_localLift_of_log_ordFrac_eq_zero}
  Proved directly in Lean by unfolding \(\ord_Q\) and applying
  \cref{lem:lineBundleAtClosedPoint_localLift_of_log_ordFrac_eq_zero} to the
  discrete valuation ring stalk at \(Q\).
\end{proof}

\section{Existence of a non-constant rational function}

\begin{corollary}



\leanok
[A genus-\(0\) curve admits a non-constant rational function with a
single simple pole]
  \label{cor:lineBundleAtClosedPoint_exists_nonconstant_genusZero}
  % Alias label consumed under the shorter name by
  % \texttt{RiemannRoch\_RationalCurveIso.tex} (RR.4):
  \label{cor:nonconstant_function_genus_zero}
  \lean{AlgebraicGeometry.Scheme.lineBundleAtClosedPoint.exists_nonconstant_genusZero}
  \uses{thm:lineBundleAtClosedPoint_dim_eq_two_of_genusZero,
        thm:lineBundleAtClosedPoint_exists_nonconstant_rational_from_dim_eq_two,
        lem:lineBundleAtClosedPoint_globalSections_iff,
        lem:lineBundleAtClosedPoint_toFunctionField_injective,
        def:principal_divisor}

  % SOURCE: [Hartshorne], IV.1, Exercise~1.1 (the existence of a
  % nonconstant rational function regular everywhere except at one
  % point on a genus-\(0\) curve, as a direct consequence of the
  % Riemann--Roch dimension formula) (read from
  % references/hartshorne-algebraic-geometry.pdf, PDF page 314).
  % SOURCE QUOTE: "1.1. Let \(X\) be a curve, and let \(P \in X\) be a
  % point. Then there exists a nonconstant rational function $f \in
  % K(X)\(, which is regular everywhere except at \)P$."
  \textit{Source: Hartshorne, IV.1, Exercise~1.1, p.~297.}
  Let \(C\) be a smooth proper geometrically irreducible curve over
  \(\bar k\) with \(g(C) = 0\), and let \(P \in C\) be a closed point. Then
  there exists a non-constant rational function \(f \in K(C)\) such that:
  \begin{itemize}
    \item \(f\) is regular on \(C \setminus \{P\}\), i.e.\ \(\ord_Q(f) \geq 0\)
      for every closed point \(Q \in C\) with \(Q \neq P\);
    \item \(f\) has at most a simple pole at \(P\), i.e.\
      \(\ord_P(f) \geq -1\);
    \item \(f \notin \bar k\), equivalently
      \(\div(f) \neq 0 \in \mathrm{Div}(C)\).
  \end{itemize}
  In particular, the pair \((1, f)\) spans
  \(H^0(C, \mathcal O_C(P))\) as a \(\bar k\)-vector space.
\end{corollary}

\begin{proof}
  
  
  
  By
  \cref{thm:lineBundleAtClosedPoint_dim_eq_two_of_genusZero},
  \(\dim_{\bar k} H^0(C, \mathcal O_C(P)) = 2\). The constant subspace
  \(\bar k \cdot 1 \subset H^0(C, \mathcal O_C(P))\), obtained from the
  structural inclusion $\bar k \cong H^0(C, \mathcal O_C)
  \hookrightarrow H^0(C, \mathcal O_C(P))$ (the rational function
  \(f = 1\), with \(\div(1) + [P] = [P] \geq 0\), under the
  identification of
  \cref{lem:lineBundleAtClosedPoint_globalSections_iff}), is a one-dimensional
  \(\bar k\)-vector subspace. Since the ambient space has dimension \(2\),
  the quotient \(H^0(C, \mathcal O_C(P)) / \bar k \cdot 1\) is a
  non-zero one-dimensional \(\bar k\)-vector space; any lift to
  \(H^0(C, \mathcal O_C(P))\) of a non-zero element of the quotient is
  an element of \(H^0(C, \mathcal O_C(P)) \setminus \bar k \cdot 1\).

  Choose such a lift $s \in H^0(C, \mathcal O_C(P)) \setminus \bar k
  \cdot 1$. Under the identification of
  \cref{lem:lineBundleAtClosedPoint_globalSections_iff} the section
  \(s\) corresponds to a rational function \(f \in K(C)^{\times}\)
  satisfying \(\div(f) + [P] \geq 0\), equivalently \(\ord_Q(f) \geq 0\)
  for every \(Q \neq P\) and \(\ord_P(f) \geq -1\). The choice
  \(s \notin \bar k \cdot 1\) means \(f\) is not the image of a constant,
  so \(f \notin \bar k\) (every non-zero element of
  \(H^0(C, \mathcal O_C) \cong \bar k\) corresponds, under the
  identification, to a constant element of \(K(C)\)). The three bullets
  of the statement are now immediate.

  Spanning. Linear independence of \(1\) and \(f\) over \(\bar k\) in
  \(H^0(C, \mathcal O_C(P))\) is the same condition as
  \(s \notin \bar k \cdot 1\) above; combined with the dimension count
  \(\dim_{\bar k} H^0(C, \mathcal O_C(P)) = 2\), the pair \((1, f)\) is a
  \(\bar k\)-basis of \(H^0(C, \mathcal O_C(P))\).

  \paragraph{Lean signature scope.} The Lean signature
  \texttt{AlgebraicGeometry.Scheme.lineBundleAtClosedPoint.exists\_nonconstant\_genusZero}
  returns an element \(f\) of the project's function-field carrier
  \(K(C)\) together with proofs of regularity on \(C \setminus \{P\}\) and
  the simple-pole bound at \(P\), packaged in the same shape as
  \cref{thm:lineBundleAtClosedPoint_exists_nonconstant_rational_from_dim_eq_two}'s output. The
  non-constancy is recorded as a separate field; the spanning
  statement is downstream and is consumed by \texttt{RR.4} via
  \texttt{Proj.fromOfGlobalSections}.
\end{proof}

\section{Mathlib-status flag: the H\(^1\)-vanishing input}

The vanishing \(H^1(C, \mathcal O_C(P)) = 0\) on a genus-\(0\) curve
(\cref{lem:H1_vanishing_lineBundleAtClosedPoint_genusZero}) is the
critical cohomological input of this chapter. Its closure depends on
two Mathlib-status facts, both of which the project already consumes
elsewhere:

\begin{itemize}
  \item \textbf{Long exact sequence of \(\Module \bar k\)-flavoured sheaf
    cohomology.} Mathlib (\texttt{b80f227}) ships
    \texttt{CategoryTheory.Abelian.Ext.covariantSequence\_exact} for
    \(\Ext^*\)-long-exact-sequences at the abelian-group level. The
    project's \(\Module \bar k\)-flavoured pipeline
    (\cref{def:Scheme_HModule}, \texttt{Cohomology\_StructureSheafModuleK.tex})
    inherits it by forget-functor naturality. This bridge is already
    consumed by \cref{thm:euler_char_eq_deg_plus_one_minus_genus} of
    \texttt{RR.2}, so closing it here introduces no new Mathlib
    obligation.
  \item \textbf{Vanishing of higher cohomology of a skyscraper sheaf}
    on a Noetherian scheme. Hartshorne~III.2.5 gives the flasque
    formulation; the project's affine-cover Čech route
    (\cref{def:Scheme_HModule} and surrounding
    \texttt{Cohomology\_StructureSheafModuleK.tex} machinery) computes
    the same vanishing for the skyscraper at a closed point by direct
    Čech-degeneration. \texttt{RR.2}'s inductive step already invokes
    \(\chi(k(P)) = 1\), which is the combined statement
    \(\dim H^0(C, k(P)) = 1\) and \(\dim H^i(C, k(P)) = 0\) for \(i \geq 1\),
    so the vanishing \(H^1(C, k(P)) = 0\) is already a consumer
    obligation of \texttt{RR.2}.
\end{itemize}

Consequently the H\(^1\)-vanishing
\cref{lem:H1_vanishing_lineBundleAtClosedPoint_genusZero} introduces no
Mathlib obligation that is not already on \texttt{RR.2}'s critical
path. The genus-\(0\) specialisation should therefore close at the
formalisation stage without splitting this chapter or pushing the
vanishing back into \texttt{RR.2}.

% NOTE (iter-174 plan-agent flag). The general-\(g\) vanishing statement
% ``\(H^1(C, \mathcal L) = 0\) whenever \(\deg \mathcal L > 2g - 2\)''
% (Hartshorne~IV.1, the Riemann--Roch consequence read covariantly) is
% \emph{not} a deliverable of this chapter and is not required for the
% genus-\(0\) critical path: the genus-\(0\) argument closes via the long
% exact sequence at \(D = 0\) without ever needing the general bound. The
% general bound is a Serre-duality-flavoured statement and is queued for
% the post-RR-bridge follow-up chapter together with the dualizing sheaf
% \(\omega_C\) and the canonical-divisor identification \(K_C \sim -2[P]\)
% on \(\mathbb P^1\) (Hartshorne~II.8.20.1).

\section{Out of scope}

The following items are deliberately \emph{not} part of \texttt{RR.3}
and are sequenced into the sibling sub-build chapters or queued for
follow-up iterations:

\begin{itemize}
  \item The construction of the invertible sheaf \(\mathcal O_C(D)\)
    associated to an arbitrary Weil divisor $D = \sum n_i [P_i] \in
    \mathrm{Div}(C)$, and the linear-equivalence isomorphism
    \(\mathcal O_C(D) \cong \mathcal O_C(D')\) when \(D \sim D'\), are
    \emph{not} performed in this chapter. The chapter constructs only
    the specialisation \(\mathcal O_C(D)\) at \(D = [P]\) a prime divisor,
    which is the only case consumed by the genus-\(0\) classification
    chain. The general \(\mathcal O_C(D)\) is queued for the
    post-RR-bridge follow-up (where it pairs with the general-\(g\)
    Riemann--Roch formula and the canonical-divisor identification).
  \item The general-\(g\) statement
    $\dim_{\bar k} H^0(C, \mathcal O_C(P)) = \deg([P]) + 1 -
    g(C) + \dim_{\bar k} H^1(C, \mathcal O_C(P))$, while a direct
    consequence of
    \cref{thm:euler_char_eq_deg_plus_one_minus_genus} of \texttt{RR.2},
    is not separately formalised here; the chapter restricts to the
    genus-\(0\) specialisation \(g = 0\) on which
    \cref{thm:lineBundleAtClosedPoint_dim_eq_two_of_genusZero} is
    available.
  \item \textbf{Serre duality} $H^1(C, \mathcal L) \cong
    H^0(C, \omega_C \otimes \mathcal L^{-1})^*$ and the associated
    canonical divisor \(K_C = c_1(\omega_C)\) are \emph{not} used at any
    step of this chapter and are not deliverables; the chapter routes
    cohomological vanishing through the long exact sequence at
    \(D = 0\) (\cref{lem:H1_vanishing_lineBundleAtClosedPoint_genusZero})
    instead. Serre duality is queued for the post-RR-bridge follow-up.
  \item The morphism \(C \to \mathbb P^1_{\bar k}\) produced by applying
    \texttt{Proj.fromOfGlobalSections} to the basis \((1, f)\) of
    \cref{cor:lineBundleAtClosedPoint_exists_nonconstant_genusZero},
    and its identification as an isomorphism, are the content of
    \texttt{RR.4} (sibling chapter
    \texttt{RiemannRoch\_RationalIsoP1.tex}, to be added). The present
    chapter supplies the basis \((1, f)\) and stops.
  \item The general-\(g\) vanishing
    \(H^1(C, \mathcal L) = 0\) for \(\deg \mathcal L > 2g - 2\) is
    \emph{not} part of this chapter and is not required for the
    genus-\(0\) critical path. It is a Serre-duality-flavoured statement
    queued for the post-RR-bridge follow-up.
  \item The interpretation of \(H^0(C, \mathcal O_C(P))\) as the space
    of sections of a complete linear system \(|[P]|\) in
    projective-space form (Hartshorne~II.7.7) is gestured at by
    \cref{lem:lineBundleAtClosedPoint_globalSections_iff} but the
    projective-space structure on the underlying set is not separately
    formalised; the chapter consumes only the \(\bar k\)-vector space
    structure on \(H^0(C, \mathcal O_C(P))\).
\end{itemize}
